\section{Discussion and Scope}
\label{sec:discussion}

\paragraph{What a quantity-preserving reference contributes.}
Maximizing feasible count makes the opportunity cost of a secondary acquisition score observable. The scientific question begins after that arithmetic fact: does the score earn the observations it displaces? FAVE, IG, and the random-order control all produce sizeable quantity differences when the maximum-count equality is removed, yet composition-first selection has no consistent predictive advantage. Conversely, the additional labels preserved by LRBE do not reliably improve prediction; UCI Brier is slightly worse for LRBE-IG than for QualityFirst-IG. The useful outcome is therefore a reference comparison, not a universal ordering. Within the eligibility rule and predicted-risk constraint, MinRisk states the least predicted risk required to retain the current maximum count.

\paragraph{Policy effects and composition evidence.}
The three experimental views answer different questions. Dynamic replay compares complete acquisition policies that begin with the same data stream but evolve through policy-specific feedback, predictions, candidates, and budgets. Same-state audits fix those quantities and show exactly how selectors exchange count, risk, and score at one decision point. The controlled experiment adds a single retraining step after fixing the acquisition state. It reproduces the count and predicted-risk contrasts, while its exposure and target-weight interactions remain uncertain. Taken together, the results locate the trade-off at selection but do not reduce long-run policy differences to composition alone.

\paragraph{Why predictive differences can remain small.}
Most labels added in the controlled task come from ordinary approvals shared across selectors: about 3,751 per incoming batch, compared with exploratory $K^*$ means of 29.55 under unit exposure and 118.7 under heterogeneous exposure. Differences among exploratory selections therefore modify only a fraction of the post-acquisition training set. This scale contrast offers one plausible account of the small Brier differences, but the experiment was designed to test exposure and target-weight interactions rather than isolate dilution. More broadly, an acquisition score is useful only when the labels it changes affect the fitted predictor and the target metric; improved score totals or altered feedback composition are intermediate outcomes.

\paragraph{Scope of inference.}
The evidence concerns random-order replay on historical cohorts with observed outcomes. LendingClub does not contain outcomes for actual historical rejects, labels arrive immediately in the simulation, and $\tau=0.5$ is a common experimental operating point rather than a production cutoff. PPDai contributes unit-exposure exploratory evidence; monetary heterogeneity comes from HMEQ and LendingClub. The IG comparison uses the registered primary budget $\rho=0.1$, while the broader budget sweep uses the original policy set. Published online exploration methods motivate the question, but the randomized controls are not implementations of those algorithms. The conclusions apply to the evaluated acquisition policies and replay conditions.

\paragraph{Risk interpretation.}
An optimizer can enforce a sum of predicted losses without controlling the tail of realized loss. Here, an exceedance means that realized defaulted exposure in a seed--round observation is larger than the policy's point-estimate budget. It is an ex-post comparison with an expected-loss quantity, not a direct estimate of safety failure. The observed frequency does not attribute the difference separately to probability error, default randomness, or exposure assumptions. Live credit experimentation would require product-specific exposure and recovery definitions, delayed-outcome handling, concentration and fairness safeguards, and explicit stopping rules. None of these operational guarantees follows from the current point-cost constraint.

\section{Conclusion}
\label{sec:conclusion}
Risk-budgeted feedback acquisition involves two distinct choices: how many outcomes to make observable and which applicants should provide them. LRBE separates these choices by preserving the current maximum feasible count before applying a composition score. Across FAVE, independent IG, and randomized controls, removing that constraint often sacrifices substantial feedback, but preserving more labels does not itself guarantee better prediction. At nearly the same cumulative count, FAVE and IG also spend more predicted risk than MinRisk without a consistent downstream return. The controlled analysis reproduces these selection trade-offs but does not identify exposure heterogeneity or target-region emphasis as reliable moderators. Acquisition studies should therefore report quantity, predicted-risk use, composition change, and downstream performance together. Minimum-risk maximum-count selection provides the reference needed to decide whether a more elaborate score earns its cost.
